\subsection{Grammars and Parsing}
\label{sec:grammars-parsing}
\slidesrc{Syntax: sentence, rules: 7--17}

\begin{nbdef}{Parsing}
In computational linguistics, \textbf{parsing} is the formal analysis by a computer of a
sentence (or other string of words) into its \textbf{constituents}, resulting in a
\textbf{parse tree} that shows their syntactic relations to each other. The tree may also
contain semantic and other information.
\end{nbdef}

Parsing is one of the possible ways to address the POS problem: once the sentence has a
tree, every word has a category (the node right above it).

\subsubsection{Chomsky and generative grammar}
\slidesrc{Syntax: sentence, rules: 8--9}

In \emph{Syntactic Structures} (1957), Noam Chomsky proposed the theory of
\textbf{generative grammar}: a finite set of rules that recursively ``specify'' or
``generate'' all the well-formed expressions of a natural language (and only them). It
is considered one of the most significant contributions to theoretical linguistics of
the 20th century.

Recursion is what links generative grammar to the question of
Section~\ref{sec:human-vs-animal}. ``Animal languages'' are forms of animal communication
that show similarities to human language: the inventory of signs may be large, the signs
relatively arbitrary, and they seem to be produced with some volition. But researchers
agree that they are not as complex or expressive as human language: they lack the
\emph{creation of new patterns of signs under varied circumstances}, i.e.\ the
productivity that recursive rules provide.

\subsubsection{Grammars and the Chomsky hierarchy}
\slidesrc{Syntax: sentence, rules: 10--13}

\begin{nbdef}{Grammar and formal language}
A \textbf{grammar} specifies the compositional structure of complex messages. A
\textbf{formal language} is a set of strings of \textbf{terminal symbols} (the words).
Each string of the language can be analysed (recognised) or generated by the grammar,
which is a set of \textbf{rewrite rules}, e.g.
\begin{align*}
  S &\to NP\ VP\\
  Article &\to \textbf{the} \mid \textbf{a} \mid \textbf{an} \mid \dots
\end{align*}
Here $S$ is the \textbf{sentence} (start) symbol; $NP$ (noun phrase), $VP$ (verb
phrase) and $Article$ are \textbf{nonterminals}; \textbf{the}, \textbf{a}, \textbf{an}
are terminals; the bar $\mid$ separates alternatives.
\end{nbdef}

Restricting the \emph{form} of the rules gives classes of grammars of increasing power,
the \textbf{Chomsky--Schützenberger hierarchy}:

\begin{center}
\small
\begin{tabular}{@{}l l l l@{}}
\toprule
\textbf{Grammar} & \textbf{Languages} & \textbf{Automaton} & \textbf{Rules} \\
\midrule
Type-0 & recursively enumerable & Turing machine & no restrictions \\
Type-1 & context-sensitive & linear-bounded automaton & $\alpha A \beta \to \alpha\gamma\beta$ \\
Type-2 & context-free & (non-det.) pushdown automaton & $A \to \gamma$ \\
Type-3 & regular & finite-state automaton & $A \to aB$, \ $A \to a$ \\
\bottomrule
\end{tabular}
\captionof{table}{The Chomsky hierarchy. $A, B$: nonterminals; $a$: a terminal;
$\alpha, \beta, \gamma$: strings of terminals and nonterminals.}
\end{center}

\begin{center}
\begin{tikzpicture}
  \draw[fill=nbblue!6,   draw=nbblue]   (0,0.25) ellipse (4.4 and 2.3);
  \draw[fill=nbgreen!8,  draw=nbgreen]  (0,-0.35) ellipse (3.3 and 1.55);
  \draw[fill=nborange!10,draw=nborange] (0,-0.7) ellipse (2.4 and 1.1);
  \draw[fill=nbpurple!10,draw=nbpurple] (0,-1.0) ellipse (1.4 and 0.65);
  \node[font=\small] at (0,1.85) {recursively enumerable};
  \node[font=\small] at (0,0.8) {context-sensitive};
  \node[font=\small] at (0,0.0) {context-free};
  \node[font=\small] at (0,-1.0) {regular};
\end{tikzpicture}
\captionof{figure}{Each class contains the ones below it. Regular expressions
(Section~\ref{sec:regular-expressions}) describe exactly the regular languages.}
\end{center}

In words:
\begin{itemize}[nosep]
  \item \textbf{Regular}: \emph{nonterminal} $\to$ \emph{terminal} [\emph{nonterminal}],
        e.g.\ $S \to \mathbf{a}S$, $S \to \Lambda$ (the empty string): any number of
        \textbf{a}. Recognised with a fixed, finite memory.
  \item \textbf{Context-free}: \emph{nonterminal} $\to$ anything, e.g.\
        $S \to \mathbf{a}S\mathbf{b}$: generates $\mathbf{a}^n\mathbf{b}^n$, which needs
        to \emph{count} the a's, impossible with finite memory (a stack suffices).
  \item \textbf{Context-sensitive}: more nonterminals on the left-hand side, e.g.\
        $ASB \to AA\mathbf{a}BB$: $S$ can be rewritten only when surrounded by $A$ and
        $B$.
  \item \textbf{Recursively enumerable}: no constraints.
\end{itemize}

\begin{nbwarning}[``natural languages are probably context-free'']
The slides say that natural languages are probably context-free, hence parsable in real
time. This is a good working assumption: context-free grammars cover most syntactic
constructions and are the basis of most parsers. Strictly speaking, however, a few
constructions are known to go beyond context-free power, e.g.\ the cross-serial
dependencies of Swiss German (Shieber, 1985). Linguists therefore say that natural
languages are \emph{mildly} context-sensitive.
\end{nbwarning}

\subsubsection{A grammar for the Wumpus world}
\slidesrc{Syntax: sentence, rules: 14--15}

The Wumpus world is a toy game from AI textbooks: an agent explores a cave with pits and
a monster, the wumpus, and perceives stenches and breezes. A tiny fragment of English for
talking about it consists of a \textbf{lexicon} and a \textbf{grammar}.

\begin{center}
\small
\begin{tabular}{@{}r c l@{}}
\multicolumn{3}{@{}l}{\textbf{Lexicon}}\\
\midrule
\emph{Noun}        & $\to$ & stench $\mid$ breeze $\mid$ glitter $\mid$ nothing $\mid$ wumpus $\mid$ pit $\mid$ pits $\mid$ gold $\mid$ east $\mid \dots$ \\
\emph{Verb}        & $\to$ & is $\mid$ see $\mid$ smell $\mid$ shoot $\mid$ feel $\mid$ stinks $\mid$ go $\mid$ grab $\mid$ carry $\mid$ kill $\mid$ turn $\mid \dots$ \\
\emph{Adjective}   & $\to$ & right $\mid$ left $\mid$ east $\mid$ south $\mid$ back $\mid$ smelly $\mid \dots$ \\
\emph{Adverb}      & $\to$ & here $\mid$ there $\mid$ nearby $\mid$ ahead $\mid$ right $\mid$ left $\mid$ east $\mid \dots$ \\
\emph{Pronoun}     & $\to$ & me $\mid$ you $\mid$ I $\mid$ it $\mid \dots$ \\
\emph{Name}        & $\to$ & John $\mid$ Mary $\mid$ Boston $\mid \dots$ \\
\emph{Article}     & $\to$ & the $\mid$ a $\mid$ an $\mid \dots$ \\
\emph{Preposition} & $\to$ & to $\mid$ in $\mid$ on $\mid$ near $\mid \dots$ \\
\emph{RelPronoun}  & $\to$ & that $\mid$ which $\mid \dots$ \\
\end{tabular}
\end{center}

The lexicon is divided into \textbf{closed classes} (articles, pronouns, prepositions:
a small fixed set, new members are almost never added) and \textbf{open classes} (nouns,
verbs, adjectives, adverbs: new words appear all the time).

\begin{center}
\small
\begin{tabular}{@{}r c l l@{}}
\multicolumn{4}{@{}l}{\textbf{Grammar}}\\
\midrule
$S$         & $\to$ & $NP\ VP$                & I + feel a breeze \\[3pt]
$Subject$   & $\to$ & $Pronoun$               & I \\
            & $\mid$ & $Noun$                 & pits \\
            & $\mid$ & $Article\ Noun$        & the + wumpus \\
            & $\mid$ & $Article\ Adjective\ Noun$ & the + smelly + wumpus \\[3pt]
$NP$        & $\to$ & $Subject$               & the wumpus \\
            & $\mid$ & $Subject\ PP$          & the wumpus + to the east \\
            & $\mid$ & $Subject\ RelClause$   & the wumpus + that is smelly \\[3pt]
$VP$        & $\to$ & $Verb$                  & stinks \\
            & $\mid$ & $Verb\ NP$             & feel + a breeze \\
            & $\mid$ & $Verb\ Adjective$      & is + smelly \\
            & $\mid$ & $Verb\ PP$             & turn + to the east \\
            & $\mid$ & $Verb\ Adverb$         & go + ahead \\[3pt]
$PP$        & $\to$ & $Preposition\ NP$       & to + the east \\
$RelClause$ & $\to$ & $RelPronoun\ VP$        & that + is smelly \\
\end{tabular}
\end{center}

All rules have a single nonterminal on the left: this is a \textbf{context-free}
grammar. Note the recursion: $NP \to Subject\ PP$ and $PP \to Preposition\ NP$, so noun
phrases can nest without limit (\emph{the wumpus to the east of the pit near the gold
\ldots}).

\begin{nbexample}{parsing ``I feel a breeze''}
\begin{center}
\begin{tikzpicture}[level distance=0.95cm,
    level 1/.style={sibling distance=4.2cm},
    level 2/.style={sibling distance=2.2cm},
    level 3/.style={sibling distance=1.4cm},
    every node/.style={font=\small}]
  \node {$S$}
    child {node {$NP$}
      child {node {$Subject$}
        child {node {$Pronoun$}
          child {node[nbblue] {I}}}}}
    child {node {$VP$}
      child {node {$Verb$}
        child {node[nbblue] {feel}}}
      child {node {$NP$}
        child {node {$Subject$}
          child {node {$Article$} child {node[nbblue] {a}}}
          child {node {$Noun$} child {node[nbblue] {breeze}}}}}};
\end{tikzpicture}
\captionof{figure}{Parse tree. The node right above each word is its part of speech:
\emph{I}/Pronoun, \emph{feel}/Verb, \emph{a}/Article, \emph{breeze}/Noun.}
\end{center}
Derivation, always rewriting the leftmost nonterminal:
$S \Rightarrow NP\ VP \Rightarrow Subject\ VP \Rightarrow Pronoun\ VP \Rightarrow
\text{I}\ VP \Rightarrow \text{I}\ Verb\ NP \Rightarrow \text{I feel}\ NP \Rightarrow
\text{I feel}\ Subject \Rightarrow \text{I feel}\ Article\ Noun \Rightarrow
\text{I feel a breeze}$.
\end{nbexample}

\begin{nbwarning}[this grammar overgenerates]
The same rules also generate \emph{*I stinks}, \emph{*me feel a breeze} and \emph{*the
wumpus feel a breeze}: nothing checks that subject and verb agree in person and number,
or that a pronoun in subject position is \emph{I} and not \emph{me}. Adding a separate
nonterminal for every combination (singular/plural, subject/object, \ldots) would make the
grammar explode. Definite Clause Grammars solve this with \emph{arguments}.
\end{nbwarning}

\subsubsection{Definite Clause Grammars}
\slidesrc{Syntax: sentence, rules: 16--17}

\begin{nbdef}{Definite Clause Grammars (DCG)}
\textbf{DCGs} are a notation for grammars natively integrated in the Prolog programming
language.
\begin{itemize}[nosep]
  \item They represent context-free grammars\ldots
  \item \ldots but nonterminals can have \textbf{arguments} (compare with \emph{attribute
        grammars}), which can express context-dependent constraints: this makes the
        formalism extremely powerful.
  \item Being Prolog, the same grammar works both as a \textbf{recogniser} (is this
        sentence grammatical?) and as a \textbf{generator} (list all grammatical
        sentences), and arbitrary reasoning can be added on top of it (``but, to your
        relief, we will not!'').
\end{itemize}
\end{nbdef}

\begin{nbexample}{agreement with an argument}
\begin{minted}{prolog}
s          --> np(Num), vp(Num).          % subject and verb share Num
np(Num)    --> article, noun(Num).
vp(Num)    --> verb(Num).
article    --> [the].
noun(sing) --> [wumpus].     noun(plur) --> [wumpuses].
verb(sing) --> [stinks].     verb(plur) --> [stink].
\end{minted}
The variable \texttt{Num} is the same in \texttt{np(Num)} and \texttt{vp(Num)}, so the
query \texttt{phrase(s, [the, wumpus, stinks])} succeeds while \texttt{phrase(s, [the,
wumpus, stink])} fails. The query \texttt{phrase(s, X)} enumerates all (two) sentences of
the language. DCGs can be run online in SWISH, \url{https://swish.swi-prolog.org/}.
\end{nbexample}

\begin{nbintuition}[why arguments add power]
Without arguments, enforcing agreement needs one copy of each rule per combination of
features ($NP_{sing}$, $NP_{plur}$, \ldots). With arguments, one rule says ``whatever the
number is, it must be the same here and there''. Arguments can carry any Prolog term,
including counters, so DCGs can describe even languages that are not context-free, such
as $\mathbf{a}^n\mathbf{b}^n\mathbf{c}^n$.
\end{nbintuition}
